\documentclass{article}
\usepackage[T1]{fontenc}
\usepackage{lmodern}
\usepackage{graphicx}
\usepackage{amsmath,amssymb}
\usepackage[a4paper,margin=2cm]{geometry}
\usepackage[absolute,overlay]{textpos}
\setlength{\TPHorizModule}{1mm}
\setlength{\TPVertModule}{1mm}
\usepackage[indonesian]{babel}
\usepackage{hyperref}
\setlength{\emergencystretch}{2em}
% unit: Halaman 24

\begin{document}

% === Halaman 24 (llm) ===

\begin{itemize}
  \item Hasil iterasi ke-2:
\end{itemize}

\vspace{5mm}

\begin{flushleft}
hereTCSWPeYraWHarSseQithaYraOeaWHstatePFairaWHKrSTYhtm\\
etheKeetPeJrSmaYPassGaseiWQhiGhitQaseWGPSseHitQasaKeaT\\
tiJTPsGaraKaeTsaWHThatthattimeTWAWSQWtSwatTraPistsSJGSTr\\
seaYreatCriUeiWasGieWtiJiGCSiWtSJioeQthereQeretQsRSTWH\\
KPaGAsCStsWearSWeeBtremitFSJtheKaGaaWHaPSWYSWeWeartheS\\
therthesGaPesQereeBGeeHiWYPFharHaWHYPSssFQithaPPtheaCC\\
earaWGeSJKTrWisheHYSPHtheQeiYhtSJtheiWseGtQasOerFremar\\
AaKPeaWHtaAiWYaPPthiWYsiWtSGSWsiHeratiSWIGSTPHharHPFKP\\
ameNTCiterJSrhisSCiWiSWresCeGtiWYit
\end{flushleft}

\vspace{10mm}

\begin{itemize}
  \item Teruskan, dengan menerka kata-kata yang sudah dikenal, misalnya\\
  \quad remarA mungkin \textit{remark}, dsb
\end{itemize}

\end{document}
